\documentclass[12pt]{article}
\usepackage{mathrsfs}
\usepackage{graphicx}
\usepackage{amssymb,amsmath,amsthm}
\usepackage{epstopdf}
\usepackage{bm}
%\usepackage{feynmf}
\usepackage[pdftex,colorlinks,urlcolor=blue]{hyperref}
\pagestyle{plain}
%\usepackage{mcite}
\DeclareGraphicsRule{.tif}{png}{.png}{`convert #1 `basename #1
.tif`.png} \textwidth = 6.5 in \textheight = 9 in \oddsidemargin =
0.0 in \evensidemargin = 0.0 in \topmargin = 0.0 in \headheight =
0.0 in \headsep = 0.0 in
\parskip = 0.1in
\parindent = 0.1in
\linespread{1.0}
\def\ba{\begin{eqnarray}}
\def\ea{\end{eqnarray}}
\def\nn{\nonumber}
\def\L{\mathcal{L}}
\def\H{\mathcal{H}}
\def\M{\mathcal{M}}
\def\N{\mathcal{N}}
\def\P{\mathcal{P}}
\def\R{\mathcal{R}}
\def\O{\mathcal{O}}
\def\q{\mathbf{q}}
\def\D{\mathcal{D}}
\def\V{\mathcal{V}}
\def\:{\boldsymbol{:}}
\def\B{\boldsymbol{B}}
\def\Bbeta{\boldsymbol{\beta}}
\def\x{\mathbf{x}}
\def\u{\mathbf{u}}
\def\y{\mathbf{y}}
\def\p{\mathbf{p}}
\def\v{\mathbf{v}}
\def\n{\hat{\mathbf{n}}}
\def\dk#1{\stackrel{\circ}{#1}}
\def\ddk#1{\stackrel{\circ \circ}{#1}}
\def\dkt#1{\stackrel{\bullet}{#1}}
\def\k{\mathbf{k}}
\def\F{\mathcal{F}}
\def\Tr{\textrm{Tr}}
\def\Vol{\textrm{Vol}}
\def\slash{\makebox[0cm][l]{/}}
\long\def\symbolfootnote[#1]#2{\begingroup%
\def\thefootnote{\fnsymbol{footnote}}\footnote[#1]{#2}\endgroup}


\begin{document}




\begin{flushright}
Chris Gauthier
\today
\end{flushright}
\begin{center} 
\vglue .06in
{\Large \bf {$U(1)$ Gauge Field Around a Black Hole}}\\[.5in]

\end{center}
%+++++++++++++++++++++++++++++++++++++++++++++++++++++++++++++++++++++++++++++++
%+++++++++++++++++++++++++++++++++++++++++++++++++++++++++++++++++++++++++++++++
%+++++++++++++++++++++++++++++++++++++++++++++++++++++++++++++++++++++++++++++++
%+++++++++++++++++++++++++++++++++++++++++++++++++++++++++++++++++++++++++++++++
\section{Quantum Corrections to Maxwell's Equations}
%+++++++++++++++++++++++++++++++++++++++++++++++++++++++++++++++++++++++++++++++
%+++++++++++++++++++++++++++++++++++++++++++++++++++++++++++++++++++++++++++++++
%+++++++++++++++++++++++++++++++++++++++++++++++++++++++++++++++++++++++++++++++
%+++++++++++++++++++++++++++++++++++++++++++++++++++++++++++++++++++++++++++++++
\subsection{Summary of the Relevant Equations}
%+++++++++++++++++++++++++++++++++++++++++++++++++++++++++++++++++++++++++++++++
%+++++++++++++++++++++++++++++++++++++++++++++++++++++++++++++++++++++++++++++++
%+++++++++++++++++++++++++++++++++++++++++++++++++++++++++++++++++++++++++++++++
First Solution:
%+++++++++++++++++++++++++++++++++++++++++++++++++++++++++++++++++++++++++++++++
\begin{gather}
A_{\mu}
=
\left\{
0,0,
\partial_{z}\Phi_{1}
,
-\partial_{\bar{z}}
\Phi_{1}
\right\}
\end{gather}
%+++++++++++++++++++++++++++++++++++++++++++++++++++++++++++++++++++++++++++++++
where
%+++++++++++++++++++++++++++++++++++++++++++++++++++++++++++++++++++++++++++++++
\begin{gather}
\partial^{2}_{t}
\Phi_{1}
-
\frac{1 - \frac{2 M}{r}}{1 + \frac{2 \xi^{2} M}{r^{3}}}
\partial_{r}
\left(
(1 - \frac{2 M}{r})
(1 + \frac{2 \xi^{2} M}{r^{3}})
\partial_{r}
\Phi_{1}
\right)
=
-
\frac{l (l + 1)}{r^{2}}
(1 - \frac{2 M}{r})
\frac{1 - \frac{4  \xi^{2} M}{r^{3}} }{1 + \frac{2 \xi^{2} M}{r^{3}}}
\Phi_{1}
\end{gather}
%+++++++++++++++++++++++++++++++++++++++++++++++++++++++++++++++++++++++++++++++
The second solution:
%+++++++++++++++++++++++++++++++++++++++++++++++++++++++++++++++++++++++++++++++
\begin{gather}
A_{\mu}
=
\left\{
0,\frac{l (l + 1) \Phi_{2}}{r^{2} (1 - \frac{4 \xi^{2} M}{r^{3}})}
,
\frac{1 - \frac{2 M}{r}}{1 + \frac{2 \xi^{2} M}{r^{3}} }
\partial_{z}
\partial_{r}
\Phi_{2}
,
\frac{1 - \frac{2 M}{r}}{1 + \frac{2 \xi^{2} M}{r^{3}} }
\partial_{\bar{z}}
\partial_{r}
\Phi_{2}
\right\}
\end{gather}
%+++++++++++++++++++++++++++++++++++++++++++++++++++++++++++++++++++++++++++++++
where
%+++++++++++++++++++++++++++++++++++++++++++++++++++++++++++++++++++++++++++++++
\begin{gather}
\partial_{t}^{2}
\Phi_{2}
-
(1 - \frac{2 M}{r})
(1 + \frac{2 \xi^{2} M}{r^{3}})
\partial_{r}
\left(
\frac{1 - \frac{2 M}{r}}{1 + \frac{2 \xi^{2} M}{r^{3}}}
\partial_{r}
\Phi_{2}
\right)
=-
\frac{l (l + 1)}{r^{2}}
(1 - \frac{2 M}{r})
\frac{1 + \frac{2 \xi^{2} M}{r^{3}}}{1 - \frac{4 \xi^{2} M}{r^{3}}}
\Phi_{2}
\end{gather}
%+++++++++++++++++++++++++++++++++++++++++++++++++++++++++++++++++++++++++++++++

%+++++++++++++++++++++++++++++++++++++++++++++++++++++++++++++++++++++++++++++++














%+++++++++++++++++++++++++++++++++++++++++++++++++++++++++++++++++++++++++++++++
%+++++++++++++++++++++++++++++++++++++++++++++++++++++++++++++++++++++++++++++++
%+++++++++++++++++++++++++++++++++++++++++++++++++++++++++++++++++++++++++++++++
%+++++++++++++++++++++++++++++++++++++++++++++++++++++++++++++++++++++++++++++++
\subsection{WKB Approximation}
%+++++++++++++++++++++++++++++++++++++++++++++++++++++++++++++++++++++++++++++++
%+++++++++++++++++++++++++++++++++++++++++++++++++++++++++++++++++++++++++++++++
%+++++++++++++++++++++++++++++++++++++++++++++++++++++++++++++++++++++++++++++++
%+++++++++++++++++++++++++++++++++++++++++++++++++++++++++++++++++++++++++++++++
Let $\Phi$ be either $\Phi_{1}$ and $\Phi_{2}$  
%+++++++++++++++++++++++++++++++++++++++++++++++++++++++++++++++++++++++++++++++
\begin{gather}
\Phi(r,t)
=
A(r)
e^{i \phi(r) - i \omega t}
Y_{l}^{m}(\theta, \varphi)
\end{gather}
%+++++++++++++++++++++++++++++++++++++++++++++++++++++++++++++++++++++++++++++++
The general structure of the kinds of equations we want to solve is
%+++++++++++++++++++++++++++++++++++++++++++++++++++++++++++++++++++++++++++++++
\begin{gather}
\partial_{t}^{2}
\Phi
-
q(r)
\frac{\partial}{\partial r}
\left(
p(r)
\frac{\partial \Phi}{\partial r}
\right)
=
-
\frac{l (l +1)}{r^{2}}
q(r) H(r)
\Phi
\end{gather}
%+++++++++++++++++++++++++++++++++++++++++++++++++++++++++++++++++++++++++++++++
In the case of $\Phi_{1}$
%+++++++++++++++++++++++++++++++++++++++++++++++++++++++++++++++++++++++++++++++
\begin{gather}
q(r)
=
\frac{1 - \frac{2 M}{r}}{1 + \frac{2 \xi^{2} M}{r^{3}}}
,
\quad
p(r)
=
(1 - \frac{2 M}{r})
(1 + \frac{2 \xi^{2} M}{r^{3}})
,
\quad
H(r)
=
1
-
\frac{4 \xi^{2} M}{r^{3}}
\end{gather}
%+++++++++++++++++++++++++++++++++++++++++++++++++++++++++++++++++++++++++++++++
In the case of $\Phi_{2}$
%+++++++++++++++++++++++++++++++++++++++++++++++++++++++++++++++++++++++++++++++
\begin{gather}
q(r)
=
(1 - \frac{2 M}{r})
(1 + \frac{2 \xi^{2} M}{r^{3}})
,
\quad
p(r)
=
\frac{1 - \frac{2 M}{r}}{1 + \frac{2 \xi^{2} M}{r^{3}}}
,
\quad
H(r)
=
\frac{1}{1 - \frac{4 \xi^{2} M}{r^{3}}}
\end{gather}
%+++++++++++++++++++++++++++++++++++++++++++++++++++++++++++++++++++++++++++++++
No matter which equation we are talking about, if we substitute the ansatz into the general equation what we get is
%+++++++++++++++++++++++++++++++++++++++++++++++++++++++++++++++++++++++++++++++
\begin{gather}
-
\omega^{2}
A
e^{i\phi}
-
q
\partial_{r}
\left[
p
(A' + i A \phi')
e^{i \phi}
\right]
=
-\frac{l (l +1)}{r^{2}}
q H
A
e^{i \phi}
\nn \\
\Rightarrow
\quad
\omega^{2}
A
+
q
\left[
p'
(A' + i A \phi')
+
p
(A'' + i A' \phi' + i A \phi'')
+
i p
(A' + i A \phi')
\phi'
\right]
=
\frac{l (l +1)}{r^{2}}
q H
A
\end{gather}
%+++++++++++++++++++++++++++++++++++++++++++++++++++++++++++++++++++++++++++++++
Assuming that $A$ and $\phi$ are real functions
%+++++++++++++++++++++++++++++++++++++++++++++++++++++++++++++++++++++++++++++++
\begin{gather}
A''
-
A \phi^{\prime 2}
+
\frac{p'}{p} A'
=
\frac{l (l +1)}{r^{2}}
\frac{H}{p} A
-
\frac{\omega^{2}}{q p}
A
\\
A \phi''
+
2 A' \phi'
+
\frac{p'}{p} A \phi'
=0
\end{gather}
%+++++++++++++++++++++++++++++++++++++++++++++++++++++++++++++++++++++++++++++++

%+++++++++++++++++++++++++++++++++++++++++++++++++++++++++++++++++++++++++++++++
\begin{gather}
\frac{1}{p}
\frac{d}{dr}
(p A')
-
A \phi^{\prime 2}
=
\left[
\frac{l (l +1)}{r^{2}}
\frac{H}{p}
-
\frac{\omega^{2}}{q p}
\right]
A
,
\quad
\frac{d}{dr}
(p A^{2} \phi')
=0
\end{gather}
%+++++++++++++++++++++++++++++++++++++++++++++++++++++++++++++++++++++++++++++++
If $\frac{d}{dr}( p A') \approx 0$ then
%+++++++++++++++++++++++++++++++++++++++++++++++++++++++++++++++++++++++++++++++
\begin{gather}
p A^{2} \phi'
=
C^{2}
\end{gather}
%+++++++++++++++++++++++++++++++++++++++++++++++++++++++++++++++++++++++++++++++
where $C$ is a constant
%+++++++++++++++++++++++++++++++++++++++++++++++++++++++++++++++++++++++++++++++
\begin{gather}
-
A
\left(
\frac{C^{2}}{p A^{2}}
\right)^{2}
=
\left[
\frac{l (l +1)}{r^{2}}
\frac{H}{p}
-
\frac{\omega^{2}}{q p}
\right]
A
\quad
\Rightarrow
\quad
\frac{C^{4}}{p}
=
\left[
\frac{\omega^{2}}{q}
-
\frac{l (l +1)}{r^{2}}
H
\right]
A^{4}
\nn \\
\Rightarrow
\quad
A
=
C
\left(
\frac{q/p}{\omega^{2}
-
\frac{l (l +1)}{r^{2}}
q H}
\right)^{1/4}
\end{gather}
%+++++++++++++++++++++++++++++++++++++++++++++++++++++++++++++++++++++++++++++++

%+++++++++++++++++++++++++++++++++++++++++++++++++++++++++++++++++++++++++++++++
\begin{gather}
\phi'
=
\frac{C^{2}}{p A^{2}}
=
\frac{1}{p}
\sqrt{
\frac{p}{q}
\left[
\omega^{2}
-
\frac{l (l +1)}{r^{2}}
q H
\right]}
=
\frac{1}{\sqrt{q p}}
\sqrt{
\omega^{2}
-
\frac{l (l +1)}{r^{2}}
q H}
\end{gather}
%+++++++++++++++++++++++++++++++++++++++++++++++++++++++++++++++++++++++++++++++
Putting this all together, for the $\Phi = \Phi_{1}$ case
%+++++++++++++++++++++++++++++++++++++++++++++++++++++++++++++++++++++++++++++++
\begin{gather}
\Phi_{1}
=
C_{\omega l m }^{(1)}
e^{- i \omega t}
Y_{l}^{m}(\theta, \varphi)
\left(
\frac{q/p}{\omega^{2}
-
\frac{l (l +1)}{r^{2}}
q H}
\right)^{1/4}
\exp\left[
\pm i \int dr \frac{1}{\sqrt{q p}}
\sqrt{
\omega^{2}
-
\frac{l (l +1)}{r^{2}}
q H} \right]
\nn \\
=
\frac{C_{\omega l m }^{(1)}
e^{- i \omega t}
Y_{l}^{m}(\theta, \varphi)}{\sqrt{1 + \frac{2 \xi^{2} M}{r^{3}}}}
\left(
\omega^{2}
-
\frac{l (l +1)}{r^{2}}
(1 - \frac{2 M}{r}) \frac{1 - \frac{4 \xi^{2} M}{r^{3}}}{1 + \frac{2 \xi^{2} M}{r^{3}}} 
\right)^{-1/4}
\nn \\
\times
\exp\left[
\pm i \int dr \frac{1}{|1 - \frac{2 M}{r}|}
\sqrt{
\omega^{2}
-
\frac{l (l +1)}{r^{2}}
(1 - \frac{2 M}{r}) \frac{1 - \frac{4 \xi^{2} M}{r^{3}}}{1 + \frac{2 \xi^{2} M}{r^{3}}} }\right]
\end{gather}
%+++++++++++++++++++++++++++++++++++++++++++++++++++++++++++++++++++++++++++++++
On the other hand, for the $\Phi = \Phi_{2}$ case
%+++++++++++++++++++++++++++++++++++++++++++++++++++++++++++++++++++++++++++++++
\begin{gather}
\Phi_{2}
=
C_{\omega l m }^{(2)}
e^{- i \omega t}
Y_{l}^{m}(\theta, \varphi)
\sqrt{1 + \frac{2 \xi^{2} M}{r^{3}}}
\left(
\omega^{2}
-
\frac{l (l +1)}{r^{2}}
(1 - \frac{2 M}{r}) \frac{1 + \frac{2 \xi^{2} M}{r^{3}}}{1 - \frac{4 \xi^{2} M}{r^{3}}} 
\right)^{-1/4}
\nn \\
\times
\exp\left[
\pm i \int dr \frac{1}{|1 - \frac{2 M}{r}|}
\sqrt{
\omega^{2}
-
\frac{l (l +1)}{r^{2}}
(1 - \frac{2 M}{r}) \frac{1+ \frac{2 \xi^{2} M}{r^{3}}}{1 - \frac{4 \xi^{2} M}{r^{3}}} }\right]
\end{gather}
%+++++++++++++++++++++++++++++++++++++++++++++++++++++++++++++++++++++++++++++++

%+++++++++++++++++++++++++++++++++++++++++++++++++++++++++++++++++++++++++++++++






%+++++++++++++++++++++++++++++++++++++++++++++++++++++++++++++++++++++++++++++++
%+++++++++++++++++++++++++++++++++++++++++++++++++++++++++++++++++++++++++++++++
%+++++++++++++++++++++++++++++++++++++++++++++++++++++++++++++++++++++++++++++++
\section{Transmission Amplitude}
%+++++++++++++++++++++++++++++++++++++++++++++++++++++++++++++++++++++++++++++++
%+++++++++++++++++++++++++++++++++++++++++++++++++++++++++++++++++++++++++++++++
%+++++++++++++++++++++++++++++++++++++++++++++++++++++++++++++++++++++++++++++++

%+++++++++++++++++++++++++++++++++++++++++++++++++++++++++++++++++++++++++++++++
\begin{gather}
U(r)
=
\frac{l (l + 1)}{r^{2}}
(1 - \frac{2 M}{r})
\end{gather}
%+++++++++++++++++++++++++++++++++++++++++++++++++++++++++++++++++++++++++++++++

%+++++++++++++++++++++++++++++++++++++++++++++++++++++++++++++++++++++++++++++++
\begin{gather}
U_{1}(r)
=
\frac{l (l + 1)}{r^{2}}
(1 - \frac{2 M}{r})
\frac{1 - \frac{4 \xi^{2} M}{r^{3}}}{1 + \frac{2 \xi^{2} M}{r^{3}}}
,
\quad
U_{2}(r)
=
\frac{l (l + 1)}{r^{2}}
(1 - \frac{2 M}{r})
\frac{1 + \frac{2 \xi^{2} M}{r^{3}}}{1 - \frac{4 \xi^{2} M}{r^{3}}}
\end{gather}
%+++++++++++++++++++++++++++++++++++++++++++++++++++++++++++++++++++++++++++++++
The transmission amplitude
%+++++++++++++++++++++++++++++++++++++++++++++++++++++++++++++++++++++++++++++++
\begin{gather}
T_{l}(\omega)
=
\frac{}{}
\end{gather}
%+++++++++++++++++++++++++++++++++++++++++++++++++++++++++++++++++++++++++++++++
where
%+++++++++++++++++++++++++++++++++++++++++++++++++++++++++++++++++++++++++++++++
\begin{gather}
\theta_{l}(\omega)
=
\int_{r_{1}}^{r_{2}}
\frac{dr}{1 - \frac{2 M}{r}}
\sqrt{U(r) - \omega^{2}}
\end{gather}
%+++++++++++++++++++++++++++++++++++++++++++++++++++++++++++++++++++++++++++++++
Find the maximum of $U(r)$
%+++++++++++++++++++++++++++++++++++++++++++++++++++++++++++++++++++++++++++++++
\begin{gather}
\left.
\frac{dU(r)}{dr} 
\right|_{r = r_{max}}
=
-
\frac{2 l (l + 1)}{r_{max}^{3}}
(1 - \frac{3 M}{r_{max}})
=0
\quad
\Rightarrow
\quad
r_{max}
=
3 M
\end{gather}
%+++++++++++++++++++++++++++++++++++++++++++++++++++++++++++++++++++++++++++++++
$U(r)$ at maximum
%+++++++++++++++++++++++++++++++++++++++++++++++++++++++++++++++++++++++++++++++
\begin{gather}
U_{max} = U(r_{max})
=
\frac{l (l + 1)}{9 M^{2}}
(1- \frac{2}{3})
=
\frac{l (l + 1)}{27 M^{2}}
\end{gather}
%+++++++++++++++++++++++++++++++++++++++++++++++++++++++++++++++++++++++++++++++
As $l$ increases the top of the barrier increases as well. For a given frequency $\omega$ there will be some $l_{c}$ at which $\omega^{2} < U_{max,l_{c}}$. The $l_{c}$ at which this happens is
%+++++++++++++++++++++++++++++++++++++++++++++++++++++++++++++++++++++++++++++++
\begin{gather}
\omega^{2} = U_{max,l_{c}}
=
\frac{l_{c} (l_{c} + 1)}{27 M^{2}}
\quad
\Rightarrow
\quad
l_{c} (l_{c} + 1)
=
27 \omega^{2} M^{2}
\label{lcdef}
\end{gather}
%+++++++++++++++++++++++++++++++++++++++++++++++++++++++++++++++++++++++++++++++
The expression for the absorption cross section is 
%+++++++++++++++++++++++++++++++++++++++++++++++++++++++++++++++++++++++++++++++
\begin{gather}
\sigma_{abs}
=
\frac{\pi}{2 \omega^{2}}
\sum_{l=1}^{\infty}
(2 l + 1)
\left[
2 - R_{l}^{(+)} - R_{l}^{(\times)}
\right]
=
\frac{\pi}{2 \omega^{2}}
\sum_{l=1}^{\infty}
(2 l + 1)
\left[
T_{l}^{(+)} + T_{l}^{(\times)}
\right]
\end{gather}
%+++++++++++++++++++++++++++++++++++++++++++++++++++++++++++++++++++++++++++++++
where $R_{l}^{(+)}$ and $R_{l}^{(\times)}$ are the reflection probabilities for the two polarization states, and $T_{l}^{(+)}$ and $T_{l}^{(\times)}$ are the corresponding transmission probabilities. In this case $T_{l}^{(+)} = T_{l}^{(\times)} = T_{l}$. Note that for $\omega^{2} < U_{max}$ the transmission probability is essentially equal to zero. On the other hand, if $\omega^{2} > U_{max}$ the transmission probability is virtually 100\%. So
%+++++++++++++++++++++++++++++++++++++++++++++++++++++++++++++++++++++++++++++++
\begin{gather}
T_{l} 
=
\left\{
\begin{array}{cc}
1 & \textrm{if } l < l_{c}
\\
0 & \textrm{if } l > l_{c}
\end{array}
\right.
\end{gather}
%+++++++++++++++++++++++++++++++++++++++++++++++++++++++++++++++++++++++++++++++
Thus,
%+++++++++++++++++++++++++++++++++++++++++++++++++++++++++++++++++++++++++++++++
\begin{gather}
\sigma_{abs}
=
\frac{\pi}{\omega^{2}}
\sum_{l=1}^{\infty}
(2 l + 1)
T_{l}
=
\frac{\pi}{\omega^{2}}
\sum_{l=1}^{l_{c}}
(2 l + 1)
=
\frac{\pi l_{c} (l_{c} + 2)}{\omega^{2}}
\end{gather}
%+++++++++++++++++++++++++++++++++++++++++++++++++++++++++++++++++++++++++++++++
Taking (\ref{lcdef}) 
%+++++++++++++++++++++++++++++++++++++++++++++++++++++++++++++++++++++++++++++++
\begin{gather}
l_{c} (l_{c} + 1)
=
27 \omega^{2} M^{2}
\quad
\Rightarrow
\quad
l_{c} = \frac{\sqrt{1 + 108 \omega^{2} M^{2}} - 1}{2}
\end{gather}
%+++++++++++++++++++++++++++++++++++++++++++++++++++++++++++++++++++++++++++++++
Substituting
%+++++++++++++++++++++++++++++++++++++++++++++++++++++++++++++++++++++++++++++++
\begin{gather}
\sigma_{abs}
=
\frac{\pi}{2 \omega^{2}}
(\sqrt{1 + 108 \omega^{2} M^{2}} -1 + 54 \omega^{2} M^{2})
\approx
54 \pi M^{2}
-
729
\pi \omega^{2} M^{4}
\end{gather}
%+++++++++++++++++++++++++++++++++++++++++++++++++++++++++++++++++++++++++++++++

%+++++++++++++++++++++++++++++++++++++++++++++++++++++++++++++++++++++++++++++++
\begin{gather}
\left.
\frac{dU_{1}(r)}{dr} 
\right|_{r = r_{max} + \delta r}
\nn \\
=
-
\left.\frac{2 l (l + 1)}{r^{4}}
\frac{(r^{6} - 8 \xi^{4} M^{2}) (r - 3 M) + (24 M - 11 r) r^{3} \xi^{2} M}{(r^{3} + 2 \xi^{2} M)^{2}}
\right|_{r = 3 M + \delta r}
=0
\nn \\
\Rightarrow
\quad
\left.
(r^{6} - 8 \xi^{4} M^{2}) (r - 3 M) + (24 M - 11 r) r^{3} \xi^{2} M 
\right|_{r = 3 M + \delta r}
= 0
\nn \\
\Rightarrow
\quad
- 243 \xi^{2} M^{5}
+
(729 M^{6} - 540 \xi^{2} M^{4} - 8 \xi^{4} M^{2})
\delta r
\approx 0
\nn \\
\Rightarrow
\quad
\delta r
\approx
\frac{243 \xi^{2} M^{3}}{729 M^{4} - 540 \xi^{2} M^{2} - 8 \xi^{4}}
\approx
\frac{\xi^{2}}{3 M}
\end{gather}
%+++++++++++++++++++++++++++++++++++++++++++++++++++++++++++++++++++++++++++++++

%+++++++++++++++++++++++++++++++++++++++++++++++++++++++++++++++++++++++++++++++
\begin{gather}
U_{1}(3 M + \delta r)
\approx
\frac{l (l +1)}{27 M^{2}}
-
\frac{2 l (l+1) \xi^{2}}{243 M^{4}}
=
\frac{l (l +1)}{27 M^{2}}
\left(
1
-
\frac{2 \xi^{2}}{9 M^{2}}
\right)
\end{gather}
%+++++++++++++++++++++++++++++++++++++++++++++++++++++++++++++++++++++++++++++++

%+++++++++++++++++++++++++++++++++++++++++++++++++++++++++++++++++++++++++++++++
\begin{gather}
\omega^{2}
=
\frac{l_{c} (l_{c} +1)}{27 M^{2}}
\left(
1
-
\frac{2 \xi^{2}}{9 M^{2}}
\right)
\quad
\Rightarrow
\quad
l_{c} (l_{c} + 1)
=
\frac{27 \omega^{2} M^{2}}{1- \frac{2 \xi^{2}}{9 M^{2}}}
\end{gather}
%+++++++++++++++++++++++++++++++++++++++++++++++++++++++++++++++++++++++++++++++

%+++++++++++++++++++++++++++++++++++++++++++++++++++++++++++++++++++++++++++++++
\begin{gather}
\sigma_{abs}^{(+)}
=
\frac{\pi l_{c}^{(+)}( l_{c}^{(+)} + 2)}{2 \omega^{2}}
\approx
27 \pi M^{2} + 6 \pi \xi^{2}
=
\frac{27}{4}
\pi 
r_{s}^{2}
+
6 \pi \xi^{2} 
\end{gather}
%+++++++++++++++++++++++++++++++++++++++++++++++++++++++++++++++++++++++++++++++
Similarly
%+++++++++++++++++++++++++++++++++++++++++++++++++++++++++++++++++++++++++++++++
\begin{gather}
\sigma_{abs}^{(\times)}
=
\frac{\pi l_{c}^{(\times)}( l_{c}^{(\times)} + 2)}{2 \omega^{2}}
\approx
27 \pi M^{2} - 6 \pi \xi^{2}
=
\frac{27}{4}
\pi 
r_{s}^{2}
-
6 \pi \xi^{2} 
=
\frac{27 \pi }{4}
(r_{s}^{2} - \frac{8}{9} \xi^{2})
\end{gather}




















%\bibliographystyle{h-elsevier}
%\bibliography{Neutrino-K-Essence}





\end{document}